\subsection{Second-order methods and structured linear algebra}
\label{subsec:second-order-methods}

Preconditioning rescales first-order information with a chosen metric.
Newton-type methods instead extract the metric from the Hessian itself, so each step is defined by a local quadratic model and implemented through a linear solve.
This is the point where curvature, factorization, and iterative linear algebra become inseparable.

\subsubsection{Newton's method}
\label{subsec:newton}

When $f$ is twice differentiable, the second-order Taylor approximation around $x$ is
\begin{equation}
  f(y) \approx f(x) + \nabla f(x)^\top (y-x) + \frac{1}{2}(y-x)^\top \Hess f(x) (y-x).
  \label{eq:taylor-2}
\end{equation}
If $\Hess f(x)\succ 0$, this quadratic model is strictly convex in $y$, so its unique minimizer gives the \emph{Newton step}
\begin{equation}
  \Delta x_{\mathrm{nt}} = -\Hess f(x)^{-1} \nabla f(x),
  \label{eq:newton-step}
\end{equation}
and Newton's method updates $x \leftarrow x + \Delta x_{\mathrm{nt}}$.
If $\Hess f(x)$ is singular or indefinite, the quadratic model need not have a minimizer, so the basic Newton step is not defined without modification.
If $f$ is exactly quadratic with positive definite Hessian, Newton reaches the minimizer in one step.
For the ridge logistic objective $s(\theta)$ from \cref{ex:logistic-reg}, the Newton direction is especially transparent: it solves
\[
  \bigl(X^\top W(\theta)X + mI\bigr)\Delta\theta = -\nabla s(\theta).
\]
This is the same Hessian structure that appeared earlier, now used directly to reshape the local geometry before taking a step.

\begin{figure}[t]
  \centering
  \begin{tikzpicture}
  % Choose an illustrative (non-quadratic) objective and a current iterate x_k.
  \pgfmathsetmacro{\xk}{-0.6}
  \pgfmathsetmacro{\fk}{(\xk)^4 + (\xk)^2 + 1}
  \pgfmathsetmacro{\gk}{4*(\xk)^3 + 2*(\xk)}
  \pgfmathsetmacro{\hk}{12*(\xk)^2 + 2}
  \pgfmathsetmacro{\xkp}{\xk - \gk/\hk} % minimizer of the quadratic model (full Newton step)
  \pgfmathsetmacro{\qkp}{\fk + \gk*(\xkp-\xk) + 0.5*\hk*(\xkp-\xk)^2}
  \pgfmathsetmacro{\fkp}{(\xkp)^4 + (\xkp)^2 + 1}

  \begin{axis}[
    width=0.9\linewidth,
    height=5.6cm,
    xmin=-0.95, xmax=0.9,
    ymin=0.95, ymax=3.0,
    axis lines=left,
    ticks=none,
    xlabel={$x$},
    ylabel={$f(x)$},
    clip=false,
  ]
    % True objective (illustrative strongly convex curve).
    \addplot[black,thick,domain=-0.95:0.9,samples=240] {x^4 + x^2 + 1.0};

    % Quadratic model around x_k (second-order Taylor approximation):
    %   q_k(x) = f(x_k) + f'(x_k)(x-x_k) + (1/2)f''(x_k)(x-x_k)^2.
    \addplot[stat221Red,thick,domain=-0.95:0.35,samples=240]
      {\fk + \gk*(x-\xk) + 0.5*\hk*(x-\xk)^2};

    % Mark x_k on f and the quadratic-model minimizer x_{k+1} on both q_k and f.
    \addplot[only marks,mark=*,mark size=1.6pt] coordinates {(\xk,\fk)};
    \addplot[only marks,mark=*,mark size=1.8pt,stat221Red] coordinates {(\xkp,\qkp)};

    \addplot[densely dashed,stat221Gray] coordinates {(\xk,0.95) (\xk,\fk)};
    \addplot[densely dashed,stat221Gray] coordinates {(\xkp,0.95) (\xkp,\qkp)};

    \node[anchor=north] at (axis cs:\xk,0.95) {$x_k$};
    \node[anchor=north] at (axis cs:\xkp,0.95) {$x_{k+1}$};

    \draw[-{Latex[length=3mm]},stat221Gray,thick]
      (axis cs:\xk,1.05) -- (axis cs:\xkp,1.05);

    \node[black,anchor=west] at (axis cs:0.72,2.2) {$f$};
    \node[stat221Red,anchor=west] at (axis cs:0.10,2.55) {$q_k$ (quadratic model)};
  \end{axis}
\end{tikzpicture}

  \caption{Newton's method: minimize the local quadratic (second-order Taylor) model of $f$ around $x_k$.
  The minimizer of the quadratic model gives the Newton step $x_{k+1}=x_k-\Hess f(x_k)^{-1}\nabla f(x_k)$.}
  \label{fig:newton-quadratic-model}
\end{figure}

\begin{lemma}[Newton direction is a descent direction]
If $\Hess f(x) \succ 0$ and $\nabla f(x)\neq 0$, then the Newton direction $\Delta x_{\mathrm{nt}}=-\Hess f(x)^{-1}\nabla f(x)$ satisfies $\nabla f(x)^\top \Delta x_{\mathrm{nt}} < 0$.
\end{lemma}
\begin{proof}
Let $g=\nabla f(x)$ and $H=\Hess f(x)\succ 0$.
Then
\[
\nabla f(x)^\top \Delta x_{\mathrm{nt}}
= -g^\top H^{-1} g.
\]
Since $H^{-1}\succ 0$, we have $g^\top H^{-1}g>0$ whenever $g\neq 0$, so the expression is strictly negative.
\end{proof}

\begin{definition}[Newton decrement]
When $\Hess f(x)\succ 0$, the Newton decrement is
\begin{equation}
  \lambda(x)^2 = \nabla f(x)^\top \Hess f(x)^{-1}\nabla f(x).
  \label{eq:newton-decrement}
\end{equation}
Equivalently, $\lambda(x)^2 = -\nabla f(x)^\top \Delta x_{\mathrm{nt}}$.
A common stopping criterion is $\lambda(x)^2/2 \le \varepsilon$.
\end{definition}

The basic algorithm is therefore stated for the regime where the Hessian stays positive definite along the iterates.

\begin{algorithm}
  \caption{Newton's method with backtracking line search}
  \label{alg:newton-backtracking}
  \begin{algorithmic}[1]
    \Require Twice differentiable $f$, initial point $x^{(0)}$ such that $\Hess f(x^{(k)})\succ 0$ for every iterate, tolerance $\varepsilon>0$, parameters $\alpha \in (0,1/2)$, $\beta \in (0,1)$
    \Ensure Approximate minimizer $x$
    \State $k \gets 0$
    \While{$\text{true}$}
      \State $g^{(k)} \gets \nabla f(x^{(k)})$
      \State $H^{(k)} \gets \Hess f(x^{(k)})$
      \State Solve $H^{(k)}\Delta x^{(k)} = -g^{(k)}$ for $\Delta x^{(k)}$ \Comment{Newton direction}
      \State $\lambda_k^2 \gets -{g^{(k)}}^\top \Delta x^{(k)}$ \Comment{Newton decrement squared}
      \If{$\lambda_k^2/2 \le \varepsilon$}
        \State \textbf{break}
      \EndIf
      \State $t^{(k)} \gets$ \Call{BacktrackingLineSearch}{$f, x^{(k)}, \Delta x^{(k)}, \alpha, \beta$} \Comment{\cref{alg:backtracking}}
      \State $x^{(k+1)} \gets x^{(k)} + t^{(k)}\Delta x^{(k)}$
      \State $k \gets k + 1$
    \EndWhile
    \State \Return $x^{(k)}$
  \end{algorithmic}
\end{algorithm}

The Newton decrement $\lambda(x)$ measures how much quadratic curvature would help you improve the objective at $x$.
When $\lambda(x)$ is small, the Newton step is small and the quadratic approximation is typically accurate, so Newton's method often takes full steps ($t=1$) and converges very quickly.
When $\lambda(x)$ is large, damping/backtracking is needed to stay in the region where the local quadratic model is trustworthy (see \citet[\S9.5.3]{boyd_vandenberghe_2004_convex}).

\citet[\S9.5.3]{boyd_vandenberghe_2004_convex} analyzes Newton's method in two phases:
far from the optimum, backtracking may be needed and the decrease is more modest; once close enough that the quadratic model is accurate, the method typically accepts $t=1$ and exhibits quadratic convergence.
Per-iteration cost can be high (solving linear systems in dimension $d$), and numerical issues can arise in high dimension or with ill-conditioned Hessians.

One also sometimes uses a \emph{regularized} (or damped) Newton step, solving $(\Hess f(x)+\lambda I)\Delta x = -\nabla f(x)$ with $\lambda>0$, to improve numerical stability or to ensure the system is well-conditioned.
This is already a geometric decision, not merely a linear-algebra safeguard: adding $\lambda I$ floors the local curvature and prevents the update from overreacting to directions that are poorly resolved.
When exact Hessians are unavailable or too expensive to form, the next compromise is to approximate the same local geometry rather than abandon it.

\subsubsection{Quasi-Newton updates}
\label{subsubsec:quasi-newton}

Quasi-Newton methods keep the local-model viewpoint but replace the exact inverse Hessian by a positive-definite approximation updated from gradients and steps.
Methods such as BFGS and L-BFGS are designed to preserve the secant information revealed by successive iterates while avoiding the cost of forming or factoring a full Hessian at every step; see \citet[\S6.1--\S6.2]{nocedal_wright_2006_numerical}.
In practice this often gives a useful middle ground between first-order robustness and full Newton steps: the method still follows a curvature-adapted direction, but the per-iteration linear algebra is much cheaper.

Whether the curvature comes from a true Hessian, a quasi-Newton update, or a covariance surrogate, the implementation still depends on stable positive-definite representations and efficient linear solves.
The next destinations make that computational geometry explicit: first the reusable SPD toolkit, then the complementary large-scale solver viewpoint provided by conjugate gradient.
